\documentclass[12pt]{article}
\usepackage[utf8]{inputenc}
\usepackage[T1]{fontenc}
\usepackage{hyperref}
\title{Wahgi Language Resources for K-12 Curriculum Development}
\author{Compiled by Claude Sonnet 5.5 (claude-sonnet-5-5) for Andr\'as Kornai}
\date{}
\begin{document}
\maketitle

\section{Introduction}
Wahgi (ISO 639-3 wgi; Glottocode nucl1620) is a Trans-New Guinea language of the Western Highlands of Papua New Guinea. It has about 86,000 speakers. Glottolog and WALS list North Wahgi and Middle Wahgi varieties, and Banz is an alternative name. Only a few resources turned up in the search, and they are mostly older linguistic and missionary works.

\section{Primary Educational Resources}
No primers, textbooks or school materials for Wahgi were found in the search results.

\section{Linguistic and Reference Materials}
Evelyn Ramsey (Ramsay) compiled the \emph{Middle Wahgi Dictionary} (1975), printed by SIL in Ukarumpa. It is out of print. PARADISEC has a PDF with partial, uncorrected OCR: \url{https://catalog.paradisec.org.au/collections/ER1/items/dictionary}. The collection is at \url{https://catalog.paradisec.org.au/repository/ER1}. A bibliographic record is at \url{https://glottolog.org/resource/reference/id/13029}.

Louis Luzbetak, \emph{The Middle Wahgi Dialects, Volume 1: Banz Grammar} (Catholic Mission, Banz, 1954), is recorded at \url{https://glottolog.org/resource/reference/id/56275}. Only the bibliographic record was seen.

Heinrich Aufenanger (1953) is also recorded in Glottolog at \url{https://glottolog.org/resource/reference/id/41156}.

WALS lists further sources on Wahgi phonology and morphology, including Luzbetak (1956) and Phillips (1976): \url{https://wals.info/languoid/lect/wals_code_wah}.

\section{Religious and Community Resources}
No Wahgi Bible translation or scripture portions turned up in this search. This does not show that none exist. Further checking is needed.

\section{Digital and Audio Resources}
PARADISEC (listed above) holds the digitised dictionary. No audio, video or software resources were found.

\section{Curriculum Development Context}
The language is not known to be taught in schools, and a curriculum would have to be built almost from scratch. The Ramsey dictionary and the older grammars are the basis. Further searches are needed for scripture, literacy materials, and PNG vernacular-education resources. Local partners, such as Wahgi-speaking educators and churches, would need to be identified.
\end{document}